\section{Introduction}

\subsection{Background of the Study}\label{background-of-the-study}

Urban drainage networks are essential infrastructure for managing
stormwater and reducing rainfall-induced flooding. Because drainage
components are physically connected and flow is directional, the
condition of one component can have consequences beyond the location
where the disturbance is first observed. Studies on drainage
malfunction, clogging, and blockage show that impaired components can
reduce effective conveyance, alter flood behavior, and contribute to
wider system consequences rather than remaining purely local events
\citep{haghbin2026}, \citep{funke2025,iqbal2024,wang2025}. This makes drainage assessment a network
problem rather than only a collection of isolated local inspections.

Previous drainage studies already demonstrate the practical value of
prioritization. Complex-network and graph-based methods have been used
to identify critical pipes, vulnerable components, weak links, and
high-risk pathways by combining network structure with information such
as hydraulic relevance, runoff, capacity, storage, discharge, failure
combinations, or siltation condition \citep{simone2022,simone2023,reyes2020,dastgir2023,dastgir2024,meijer2023,zhang2025}. At the same
time, computer-vision and monitoring approaches can identify local
blockage, debris, or component condition at individual assets
\citep{vandaele2024,iqbal2021,patil2023}. These two research directions address different
parts of the decision problem. Local monitoring can show how impaired a
component is at the point of observation, while graph-based approaches
can show where a component is structurally or functionally important.
The literature therefore supports the distinction between local
condition information and network-context prioritization rather than
treating them as interchangeable outputs \citep{simone2022}, \citep{dastgir2023},
\citep{meijer2023,zhang2025,vandaele2024,iqbal2021,patil2023}.

This distinction is important for inspection and maintenance decisions.
In a network containing many connected components, decision-makers may
need to determine which locations deserve earlier inspection,
maintenance, or higher-fidelity analysis. Drainage studies on
critical-pipe identification, weakest-link analysis, siltation chains,
and multiple-pipe failure show that component importance can change when
flow paths, capacity, network structure, or interacting failures are
considered \citep{dastgir2023,dastgir2024,meijer2023,zhang2025}. Consequently, two components with similar
local disturbance may differ in network-level importance, while a
structurally central component may not necessarily carry the highest
current modeled risk under the present disturbance and operating
conditions. The practical need addressed by this study is therefore not
simply to detect blockage or rank structural importance, but to compare
components according to a risk state that incorporates both local and
propagated network context.

The broader propagation literature shows that this modeling issue
extends beyond drainage. Influence-diffusion models propagate
activation, epidemic models propagate compartment or infection states,
cascading-failure models propagate load or failure, attack-graph and
cyber-risk approaches propagate reachable compromise or risk, and
financial and infrastructure models propagate exposure, loss, stress, or
related system states \citep{kempe2003,kermack1927,pastor2001,humphries2021,buldyrev2010,he2019,shuang2017,sheyner2002,munoz2017,kavallieratos2021,faraji2026,zhou2023,wu2026,petrone2018}. These model families
demonstrate that network propagation is well established, but they also
show that the quantity being propagated differs substantially across
applications. The update and aggregation rules used in each model are
therefore tied to the semantics of that state rather than being
automatically transferable to a bounded comparative risk-index task
\citep{kempe2003}, \citep{buldyrev2010}, \citep{sheyner2002,munoz2017,kavallieratos2021,faraji2026,zhou2023,wu2026,petrone2018}.

Across these model families, several recurring mechanisms are
nevertheless evident. Topology establishes where interaction may occur;
transmission or influence represents how strongly an effect transfers;
susceptibility, capacity, vulnerability, or tolerance modifies how a
relation or receiving context responds; current state influences later
propagation; and incoming contributions require an aggregation rule
consistent with the meaning of the propagated quantity \citep{kempe2003},
\citep{buldyrev2010,he2019,shuang2017,sheyner2002,munoz2017,kavallieratos2021,faraji2026,zhou2023,wu2026,petrone2018}. Drainage studies reinforce the same distinction by
showing that structural measures are often augmented with hydraulic
relevance, capacity, runoff, storage, discharge, or other operating
information when functional importance is being assessed
\citep{simone2022,simone2023,reyes2020,dastgir2023,dastgir2024,meijer2023,zhang2025}. These mechanisms provide useful building blocks, but
the reviewed evidence does not clearly establish their particular
integration into a lightweight directed-graph architecture whose primary
output is a bounded continuous non-probabilistic comparative risk state.

The present study therefore proposes the \textbf{Graph-Based Bounded Risk
Propagation Model (GBBRPM)}. The term bounded describes a core property
of the model rather than a drainage-specific disturbance: the model
represents node risk as a normalized state constrained to the interval
{[}0,1{]}. The architecture is generic in formulation, while its
parameter definitions remain domain-specific. In a drainage
instantiation, for example, local disturbance may represent blockage
severity and susceptibility may be derived from load relative to
effective capacity. A different domain may use the same core only if it
can defensibly define its own nodes, directed relations, disturbance,
susceptibility, transmission, and comparative risk semantics. The study
therefore claims generic formulation rather than universal empirical
generalizability.

\subsection{Statement of the Problem}\label{statement-of-the-problem}

The study addresses the problem of \textit{comparative risk prioritization in
interconnected directed networks}, initially motivated by urban drainage.
Existing detection and monitoring approaches can identify the condition
of an individual drainage component, including visual blockage or
obstruction, but these outputs are principally local to the monitored
asset \citep{vandaele2024,iqbal2021,patil2023}. A local blockage score may therefore indicate
that a component is impaired without directly showing how that condition
should be prioritized once upstream influence, downstream connectivity,
or relation-specific operating conditions are considered.

Graph-based structural approaches provide an important second
perspective because they use network topology to identify critical or
vulnerable components. However, drainage studies show that structural
importance and functional or operating response are not equivalent.
Simone et al. and Reyes-Silva et al. show that structural centrality may
diverge from hydraulically meaningful importance, while later
critical-pipe and weakest-link studies augment graph structure with
runoff, capacity, storage, discharge, or other domain information
\citep{simone2022,simone2023,reyes2020,dastgir2023,dastgir2024,meijer2023}. Siltation-risk analysis similarly combines
centrality, pipeline risk, edge vulnerability, and propagation pathways
rather than relying on topology alone \citep{zhang2025}. Centrality can therefore
identify structurally important locations, but it does not by itself
model how a current disturbance recursively changes as it moves through
relation-specific operating conditions.

Existing propagation models address network influence from other
directions, but their state semantics are commonly designed for
different analytical questions. Influence-diffusion models use
activation states \citep{kempe2003}, epidemic models use infection or compartment
states \citep{kermack1927,pastor2001,humphries2021}, cascading-failure models use load, capacity, or
failure states \citep{buldyrev2010,he2019,shuang2017}, attack-graph approaches use reachable
compromise or probability \citep{sheyner2002}, \citep{munoz2017}, and cyber, finance,
transportation, and infrastructure models use other forms of risk,
exposure, loss, or dynamic system state \citep{kavallieratos2021,faraji2026,zhou2023,wu2026,petrone2018}. These
formulations are appropriate for their intended purposes, but they do
not directly provide the bounded continuous non-probabilistic
comparative risk state targeted in the present study.

Taken together, the literature reveals an architectural and integrative
research problem. The individual mechanisms relevant to the target task
already exist: directed reachability is established in graph models;
recursive current-state updating appears across contagion and cascading
processes; capacity, vulnerability, susceptibility, or related response
factors appear in infrastructure and risk models; and multiple forms of
aggregation are used according to state semantics \citep{kempe2003,kermack1927,pastor2001,humphries2021,buldyrev2010,he2019,shuang2017,sheyner2002,munoz2017,kavallieratos2021,faraji2026,zhou2023,wu2026,petrone2018}.
Drainage research additionally demonstrates the practical importance of
combining structural and operating information for prioritization
\citep{simone2022,simone2023,reyes2020,dastgir2023,dastgir2024,meijer2023,zhang2025}. Within the reviewed evidence, however, these
mechanisms are not clearly consolidated into one lightweight and
interpretable architecture specifically intended to produce a bounded
comparative risk state from local disturbance and propagated network
context.

The central problem of the study is therefore how to formulate a
directed-graph model that combines local disturbance with recursively
propagated network context so that nodes can be compared according to
their current modeled risk state, while keeping topology,
susceptibility, optional transmission, aggregation, and dependence
conceptually distinct. Consistent with this central problem, the study
addresses the following research questions:

\begin{enumerate}[label=\arabic*.,leftmargin=*]
\item How can local disturbance, current upstream risk,
relation-specific susceptibility, optional transmission, directed
reachability, and bounded multi-source aggregation be integrated into
an interpretable directed-graph architecture?
\item Does the resulting formulation preserve boundedness,
monotonicity, and expected comparative-ranking behavior under
controlled network conditions?
\item How sensitive, robust, and computationally scalable is the model
under changes in disturbance severity, susceptibility, topology,
disturbance location, and multi-source configuration?
\item Can the same bounded propagation architecture be instantiated
coherently in drainage, software-dependency, and electrical-network
cases without redefining its core update rule?
\item What evidence does each domain provide for model behavior,
structural applicability, preprocessing feasibility, and the boundary
between comparative graph risk and domain-specific outcome prediction?
\end{enumerate}

\subsection{Objectives of the Study}\label{objectives-of-the-study}

\subsubsection{General Objective}\label{general-objective}

The general objective of the study is to design and computationally
evaluate a \textbf{Graph-Based Bounded Risk Propagation Model (GBBRPM)}
for interpretable comparative node-risk prioritization in directed
networks and to examine its cross-domain applicability through drainage,
software-dependency, and electrical-network cases.

\subsubsection{Specific Objectives}\label{specific-objectives}

Specifically, the study aims to:

\begin{enumerate}[label=\arabic*.,leftmargin=*]
\item synthesize the limitations and recurring mechanisms of local
condition assessment, graph-based prioritization, and existing
propagation models that are relevant to comparative network-risk
prioritization;
\item formulate a bounded continuous non-probabilistic recursive
architecture that explicitly distinguishes local disturbance,
relation-specific susceptibility, optional transmission, directed
reachability, current-state updating, bounded multi-source aggregation,
and dependence;
\item computationally validate the generic GBBRPM architecture using
controlled and reconstructed networks in terms of boundedness,
monotonicity, ranking behavior, sensitivity, robustness, reconvergence
effects, and computational scalability;
\item instantiate the frozen architecture in drainage,
software-dependency, and electrical-network cases using traceable
topology and domain-appropriate disturbance, susceptibility, and
transmission definitions; and
\item compare the resulting evidence across domains while distinguishing
behavioral validation, structural applicability, baseline-data
verification, and independent outcome validation, and identify the
model's scope boundaries and requirements for further refinement.
\end{enumerate}

\subsection{Significance of the Study}\label{significance-of-the-study}

The study is significant because it addresses a decision-support space
between purely local or structural prioritization and high-fidelity
domain simulation. Drainage literature demonstrates both the usefulness
of lightweight graph-based screening and the added physical detail
provided by hydrodynamic and learned hydraulic models \citep{simone2022},
\citep{dastgir2023}, \citep{meijer2023}, \citep{zhang2024}. GBBRPM does not attempt to replace those
higher-fidelity approaches. Instead, it is intended to produce a bounded
comparative risk state that allows components to be ranked relative to
one another under a common network and parameterization.

For urban drainage research, the study provides a way to connect local
disturbance information with network context for preliminary
prioritization of locations that may merit inspection, maintenance, or
higher-fidelity hydraulic analysis. This positioning is consistent with
drainage studies that use graph methods for critical-pipe or
weakest-link screening while retaining hydrodynamic models for richer
physical assessment \citep{dastgir2023}, \citep{meijer2023}. The model is therefore intended
as a complementary prioritization layer rather than a hydraulic
prediction model.

For graph-based risk and propagation research, the study provides an
explicit architecture in which topology, susceptibility, transmission,
local disturbance, aggregation, and dependence are treated as
conceptually distinct mechanisms rather than being collapsed into one
undifferentiated coefficient. This separation follows recurring
distinctions observed across influence, cascade, cyber-risk, financial,
and infrastructure propagation models \citep{kempe2003}, \citep{buldyrev2010,he2019,shuang2017,sheyner2002,munoz2017,kavallieratos2021,faraji2026,zhou2023,wu2026,petrone2018}. For
the drainage application and any future domain-specific application, the
model provides a reusable core propagation structure while requiring
each instantiation to justify its own parameter definitions and any
additional refinement layer.

The study is also significant methodologically because interpretability
is retained throughout the propagation process. A high modeled risk
state at a node can be traced to the current upstream state,
relation-specific susceptibility, optional transmission, local
disturbance, and the bounded aggregation mechanism. This traceability is
especially relevant to comparative prioritization because the purpose is
not only to produce an ordering, but also to explain why one component
receives a higher modeled risk state than another.

\subsection{Scope and Limitations}\label{scope-and-limitations}

The study focuses on the design, computational validation, and
cross-domain evaluation of GBBRPM as a lightweight directed-graph
architecture for comparative bounded-risk propagation. The generic core represents a node
risk state, a local disturbance, relation-specific susceptibility,
optional transmission or influence, and an edge contribution. The risk
state and the model's bounded parameters are represented on normalized
{[}0,1{]} scales.

The evaluation has two complementary layers. The first evaluates the
generic formulation on directed acyclic graphs using one-pass topological
processing, controlled synthetic networks, diagnostic baselines,
perturbation trials, mathematical property tests, and scalability tests.
The second examines cross-domain instantiation through drainage,
software-dependency, and electrical-network cases. Drainage supplies
behavioral tests and a SWMM external-reference comparison; software
supplies real-topology structural applicability and transmission
sensitivity evidence; and electrical data supply real-topology and
baseline-preprocessing verification. Event-centered electrical outcome
validation remains pending because the available baseline measurements
cover normal operation rather than the documented switching events.

The study does not interpret the bounded risk state as a calibrated
probability of failure, blockage, or any other event. The base model
also does not represent cycles, bidirectional feedback, explicit
temporal dynamics, hydraulic backwater, surcharge, ponding, flow
redistribution, or common-cause dependence correction. These exclusions
are important because drainage and broader propagation literature show
that feedback, shared dependencies, and richer physical state
interactions may require iterative, probabilistic, or high-fidelity
formulations beyond a one-pass lightweight graph model \citep{buldyrev2010},
\citep{shuang2017}, \citep{wu2026}, \citep{petrone2018}, \citep{dastgir2024}, \citep{zhang2024}. Where heterogeneous
transmission cannot yet be empirically justified, the baseline
transmission term is kept neutral. Likewise, drainage-specific
susceptibility mappings such as load relative to effective capacity are
treated as instantiation choices rather than universal components of the
generic model.

The study claims generic formulation and demonstrated cross-domain
instantiability within the three evaluated cases, not universal empirical
generalizability. The cases test different claims and do not provide
equivalent outcome-validation strength. A domain may instantiate the core
only if it can provide defensible definitions for nodes, directed
relations, local disturbance, susceptibility, optional transmission, and
comparative risk. Domain-specific prediction, calibration, and operational
usefulness must be justified and validated separately from the generic
core.

\subsection{Conceptual Framework}\label{conceptual-framework}

The conceptual framework of the study begins with local disturbance and
directed network context. The current risk state at an upstream node is
transformed into a relation-specific contribution according to
susceptibility and, where justified, optional transmission. Multiple
incoming contributions are then combined with any local disturbance at
the receiving node using a bounded aggregation rule, producing an
updated comparative risk state. That updated state can in turn influence
later downstream relations, allowing the modeled risk state to evolve
recursively through the network. The resulting node states provide the
basis for comparative prioritization.

Figure~\ref{fig:conceptual-framework} distinguishes the domain-supplied
inputs, the generic bounded propagation core, and the comparative
decision-support output.

\begin{figure}[H]
\centering
\resizebox{\textwidth}{!}{%
\begin{tikzpicture}[
  font=\sffamily,
  box/.style={draw=gbgray, line width=.8pt, fill=white, minimum width=3.6cm, minimum height=1.3cm, align=center},
  core/.style={box, draw=gbnavy, line width=1.1pt, fill=gblight, minimum width=5.1cm, minimum height=2.6cm},
  arr/.style={-{Latex[length=2.4mm]}, line width=.8pt, draw=gbgray},
  note/.style={font=\sffamily\footnotesize, text=gbgray, align=center}
]
\node[box] (b) {Local disturbance\\$B_i\in[0,1]$};
\node[box, below=8mm of b] (g) {Directed topology\\$G=(V,E)$};
\node[box, below=8mm of g] (s) {Relation parameters\\$S_{ij},\;\tau_{ij}\in[0,1]$};
\node[core, right=20mm of g] (core) {\textbf{GBBRPM propagation core}\\[2mm]
$Q_{ij}=S_{ij}\tau_{ij}R_i$\\[1mm]
$R_j=1-(1-B_j)\displaystyle\prod_{i\in P(j)}(1-Q_{ij})$};
\node[box, right=20mm of core, yshift=10mm] (risk) {Current node states\\$R_i\in[0,1]$};
\node[box, below=8mm of risk] (rank) {Comparative ranking\\descending $R_i$; ties retained};
\node[note, below=5mm of rank] (decision) {Inspection, maintenance, or\\higher-fidelity-analysis support};

% Distinct connection ports along the core's west edge
\coordinate (core-in-top) at
  ($(core.north west)!0.22!(core.south west)$);

\coordinate (core-in-mid) at
  ($(core.north west)!0.50!(core.south west)$);

\coordinate (core-in-bottom) at
  ($(core.north west)!0.78!(core.south west)$);

% Input connections
\draw[arr] (b.east) -- (core-in-top);
\draw[arr] (g.east) -- (core-in-mid);
\draw[arr] (s.east) -- (core-in-bottom);

% Horizontal core-to-output connection
\draw[arr] (core.east |- risk.west) -- (risk.west);

% Remaining output flow
\draw[arr] (risk.south) -- (rank.north);
\draw[arr] (rank.south) -- (decision.north);

\node[note, below=5mm of core] {Bounded, continuous, non-probabilistic state};
\end{tikzpicture}}
\caption{Conceptual framework of GBBRPM, distinguishing supplied domain inputs, the generic propagation core, and the comparative decision-support output.}
\label{fig:conceptual-framework}
\end{figure}

For a directed relation $i\to j$, the generic edge contribution is:

\[Q_{ij}=S_{ij}\tau_{ij}R_i\eqtag{1-1}\]

where $R_i$ is the current upstream normalized risk state, $S_{ij}$ is
relation-specific susceptibility, and $\tau_{ij}$ is optional transmission or
influence.

The receiving-node state is:

\[R_j=1-(1-B_j)\prod_{i\in P(j)}(1-Q_{ij})\eqtag{1-2}\]

where $B_j$ is the local disturbance and $P(j)$ is the set of predecessors of
node $j$. The multiplicative-complement form is used as a bounded
saturating aggregation mechanism. Although it is algebraically similar
to noisy-OR, it is not interpreted as a probabilistic union because the
inputs are bounded risk contributions rather than calibrated
probabilities.

The resulting $R_j$ is the node's current modeled risk state. It may then
serve as the upstream state for later relations, allowing risk to be
recursively updated through the network. Node ranking is therefore a
decision-support output derived from the resulting risk states rather
than the propagation mechanism itself.

\subsection{Definition of Terms}\label{definition-of-terms}
 
{\renewcommand{\arraystretch}{1.15}
\begin{longtable}{@{}p{\dimexpr0.28\linewidth-\tabcolsep\relax}p{\dimexpr0.72\linewidth-\tabcolsep\relax}@{}}
\toprule
Term & Definition \\
\midrule
\endhead
\bottomrule
\endlastfoot
 
Comparative risk prioritization. &
Ranking components according to their modeled risk states relative to one another under the same network and parameterization. \\
\addlinespace
 
GBBRPM. &
Graph-Based Bounded Risk Propagation Model, the generically formulated directed-graph architecture proposed in this study. \\
\addlinespace
 
Normalized risk state. &
A dimensionless continuous value in $[0,1]$ representing the current modeled level of network-context risk at a node. It is not a calibrated probability of an event. \\
\addlinespace
 
Local disturbance. &
A bounded input representing risk or disturbance introduced directly at a node according to the selected domain. \\
\addlinespace
 
Susceptibility. &
A bounded relation-specific factor representing how strongly a relation or receiving context responds under the evaluated condition. \\
\addlinespace
 
Transmission or influence. &
An optional bounded factor describing how strongly influence transfers along an existing directed relation. \\
\addlinespace
 
Edge contribution. &
The relation-specific contribution generated from current upstream risk, susceptibility, and optional transmission. \\
\addlinespace
 
Topology or reachability. &
The directed network structure that determines which nodes and relations can participate in propagation. \\
\addlinespace
 
Bounded aggregation. &
The multiplicative-complement mechanism used to combine local disturbance and multiple incoming contributions while preserving the $[0,1]$ state range. \\
\addlinespace
 
Domain-specific instantiation. &
The process of defining the generic model's nodes, edges, disturbance, susceptibility, transmission, and any optional refinement parameters according to a particular application domain. \\
\addlinespace
 
Generic formulation. &
A model architecture whose core propagation logic is not hard-coded to one application domain. Generic formulation does not imply that the model has been empirically evaluated across all domains. \\
 
\end{longtable}
}
